\documentclass[handout, font=12pt]{ximera}


\title{Fractions: Simplifying}   % Use xmlatex all -s to compile when SERVE refuses to work.
\author{Kelly Stady}
\license{CC: 0}         % replace with an appropriate license, or set it in xmPreamble

\begin{document}  

\begin{abstract}

\end{abstract}

\maketitle

\label{xim:FractionsSimplify}


\section*{Factors of Whole Numbers}

The \textbf{factors of a whole number} are numbers that divide into it evenly (no remainder).  The factors of a number can be used to write
the number as a product.  In the example below, the number 18 has been written as a product using the factors 2 and 3.
\[ 2 \cdot 3 \cdot 3 = 18 \]

We can also write the number 18 as a product using the factors 6 and 3 or 2 and 9.  Below is a list of all factors of 18.

\begin{center}
\textbf{Factors of 18:} 1, 2, 3, 6, 9, 18
\end{center}


\begin{problem}
List all factors of the number 20.

~ \\[3em]

\begin{multipleChoice}
\choice{$1, 2, 3, 4, 5, 10, 20$}
\choice[correct]{$1, 2, 4, 5, 10, 20$}
\choice{$1, 2, 4, 5, 10, 15, 20$}
\choice{$1, 2, 4, 5, 10, 20$}
\end{multipleChoice}

\end{problem}

~ \\[3em]

\section*{Simplifying Fractions}

A fraction is said to be in \textbf{simplest form} or \textbf{lowest terms} when the
numerator and denominator have no common factors other than 1.  The process of writing a
fraction in simplest form is called \textbf{simplifying} the fraction.


The fraction $\textstyle \frac{2}{5}$ is in simplest form because 2 and 5 have no common factors.  However, the fraction
$\textstyle \frac{4}{6}$ is not in simplest form because both 4 and 6 have a common factor of 2.

~ \\[3em]

\subsection*{Greatest Common Factor (GCF)}

The \textbf{greatest common factor} of a group of numbers is the largest number that is a factor of each number.  Equivalently, it is the largest 
number that evenly divides into all of the numbers.  One way to find the GCF is to list the factors of each number and look for the \emph{largest} 
number that appears in each list.  For example, to find the GCF of 12 and 36 we begin by listing the factors of each number.



\begin{description}  %See if this works in html
\item[] \textbf{Factors of 12:} 1, 2, 3, 4, 6, 12


\item[] \textbf{Factors of 36:} 1, 3, 4, 6, 9, 12, 18, 36
\end{description}


The \textit{common factors} of 12 and 36 are 1, 3, 4, 6, and 12.  So, the \textbf{greatest common factor} is 12.

~ \\[3em]

\begin{problem}
Find the GCF of 15 and 20.

\[ \textbf{GCF} = \answer{5} \]

\end{problem}

\begin{problem}
Find the GCF of 24 and 60.

\[ \textbf{GCF} = \answer{12} \]

\end{problem}

~ \\[3em]

To use the GCF to simplify fractions, divide the numerator and denominator by the GCF.

~ \\[3em]


\begin{example}
The fraction $\textstyle \frac{4}{12}$ can be simplified by dividing the numerator and denominator by the GCF = 4.

\begin{eqnarray*}
  \frac{4}{12} &=& \frac{4 \div 4}{12 \div 4} \\
  & & \\
   &=& \frac{1}{3}
\end{eqnarray*}

\end{example}

~ \\[3em]

We can also simplify fractions by dividing the numerator and denominator by \emph{any} common factor, however using the GCF is usually the quickest method.
If we didn't notice that the GCF was 4, in the previous example, but knew that both numbers were evenly divisible by 2, we would need two steps to simplify the fraction
$\textstyle \frac{4}{12}.$
\begin{eqnarray*}
  \frac{4}{12} &=& \frac{4 \div 2}{12 \div 2} \\
  & & \\
   &=& \frac{2}{6}  \\
   & & \\
   &=& \frac{2 \div 2}{6 \div 2} \\
    & & \\
     &=& \frac{1}{3}
\end{eqnarray*}

When the numbers in the numerator and denominator are large, it can sometimes take too long to find the GCF.  In these cases, you can use the following facts to find
a common factor and proceed from there.

\begin{description}
\item[] \textbf{2 is a factor of a number} if the last digit of the number is divisible by 2.

\item[] \textbf{3 is a factor of a number} if the sum of the digits is divisible by 3.


\item[] \textbf{5 is a factor of a number} if the last digit is 0 or 5.

\end{description}

~ \\[3em]


INSERT VIDEO of simplifying a fraction using one of these facts.


\begin{problem}
Simplify the fraction.

\[ \frac{33}{60} = \answer{11/20} \]
\end{problem}

\begin{problem}
Simplify the fraction.

\[ \frac{42}{70} = \answer{3/5} \]
\end{problem}

\end{document} 